% ID: SEL-D0046-OB-E03
% ESTADO: cotejada. Segunda pasada visual y borrador compilado revisados.
% Fuente: Selectividad/sel_2014_matematicas_aplicadas/Matematicas_CCSS_JUN.pdf
% SHA-256: df7375f8183eafb6f45e9c65083e68d528ffabe07d66717b0ccf0f020c5b064f
% Página: 2; opción B; ejercicio original 3.

Antonio va a la compra dos días de cada cinco. A lo largo del tiempo, ha observado que la fruta está de oferta la tercera parte de los días que va a la compra y la mitad de los días que no va. Elegido un día al azar:
\begin{enumerate}[label=\alph*)]
\item \textbf{(1.5 puntos)} \textquestiondown{}Cuál es la probabilidad de que la fruta esté de oferta ese día?
\item \textbf{(1 punto)} Calcule la probabilidad de que ese día Antonio vaya a la compra o la fruta esté de oferta.
\end{enumerate}
